% Dance4Life - Aprendizagem por Reforço para Envelhecimento Ativo através da Dança
% Formato LNCS - UC Aprendizagem por Reforço 2025/26
%
\documentclass[runningheads]{llncs}
%
\usepackage[T1]{fontenc}
\usepackage[utf8]{inputenc}
\usepackage[portuguese]{babel}
\usepackage{graphicx}
\usepackage{booktabs}
\usepackage{amsmath}
\usepackage{hyperref}
\usepackage{url}
%
\begin{document}
%
\title{Dance4Life: Aprendizagem por Reforço para Acompanhamento de Atividade Física no Envelhecimento Ativo}
%
\titlerunning{AR para Acompanhamento no Envelhecimento Ativo}
%
\author{
Carlos Bergueira\inst{1} \and
Diego Silva\inst{1} \and
Filipa Pereira\inst{1} \and
Rui Rodrigues\inst{1}
}
%
\authorrunning{C. Bergueira et al.}
%
\institute{Universidade do Minho, Braga, Portugal \\
\email{\{pg11605,pg59999,pg42201,pg7942\}@alunos.uminho.pt}}
%
\maketitle
%
\begin{abstract}
O comportamento sedentário constitui um fator de risco major para doenças crónicas
e declínio funcional em adultos mais velhos. A dança combina exercício físico com
envolvimento social, tornando-se uma modalidade particularmente eficaz para o
envelhecimento ativo. Apresentamos o Dance4Life, uma plataforma de ponta a ponta
que recorre à Aprendizagem por Reforço (AR) para personalizar a intensidade do
acompanhamento em atividade física baseada em dança. Um ambiente Gymnasium
personalizado modela o estado tridimensional do utilizador---nível de atividade,
fadiga física e fadiga de alarmes---sobre ações de acompanhamento discretas. Dois
algoritmos de AR profunda, \emph{Proximal Policy Optimisation} (PPO) e
\emph{Deep Q-Network} (DQN), são treinados e comparados em condições idênticas.
Ambas as políticas convergem para recompensa próxima do ótimo em menos de
50\,000 passos e atingem 100\% de taxa de sobrevivência em episódios de 96 passos.
A melhor política é exportada para ONNX e implantada num smartphone Android via
ONNX Runtime, realizando inferência em tempo real a partir de dados de sensores
vestíveis. Um backend multi-agente em Python (SPADE) trata da recolha de dados,
agrupamento de utilizadores e persistência em nuvem via Firebase. Os resultados
confirmam que a AR consegue equilibrar eficazmente a promoção da atividade física
com os riscos de fadiga física e fadiga de alarmes num cenário de acompanhamento
realista.

\keywords{Aprendizagem por reforço \and Envelhecimento ativo \and Dança \and
Sensores vestíveis \and Sistemas multi-agente \and ONNX \and Android}
\end{abstract}
%
%
\section{Introdução}
\label{sec:intro}

O envelhecimento populacional é um dos desafios definidores do século~XXI.
Em 2050, o número de pessoas com 60 ou mais anos deverá atingir 2,1 mil milhões
em todo o mundo~\cite{who2022}. A inatividade física acelera a perda de mobilidade,
aumenta o risco de doenças cardiovasculares e reduz a qualidade de vida.
A dança é crescentemente reconhecida como uma intervenção simultaneamente eficaz e
agradável para adultos mais velhos: melhora o equilíbrio, a aptidão cardiovascular
e a função cognitiva, ao mesmo tempo que promove a interação
social~\cite{hwang2016}.

Um desafio central no acompanhamento mediado por tecnologia é a
\emph{personalização}. Um sinal de acompanhamento que motiva um utilizador pode
irritar outro; uma sessão demasiado intensa pode causar danos físicos, ao passo
que uma demasiado suave pode não produzir benefício algum. As abordagens estáticas
baseadas em regras não conseguem adaptar-se ao estado fisiológico e psicológico em
constante mudança do utilizador.

Propomos o Dance4Life, uma plataforma que aborda este desafio através da AR.
O sistema aprende uma política de acompanhamento que seleciona a intensidade de
intervenção adequada em cada ponto de decisão, observando um vetor de estado
compacto derivado de sensores vestíveis e otimizando uma função de recompensa que
simultaneamente promove a atividade, penaliza o comportamento sedentário e previne
eventos críticos de fadiga física ou fadiga de alarmes.

As principais contribuições deste trabalho são:
\begin{itemize}
  \item Um ambiente Gymnasium personalizado que captura o compromisso entre
        promoção de atividade, fadiga física e irritação do utilizador num
        contexto de acompanhamento por dança.
  \item Uma avaliação comparativa de PPO e DQN em condições de treino idênticas,
        com métricas quantitativas.
  \item Um pipeline de implantação de ponta a ponta, desde a política treinada até
        à inferência ONNX no dispositivo Android.
  \item Um backend multi-agente que orquestra a recolha de dados de sensores,
        o enriquecimento contextual, o agrupamento de utilizadores e a
        persistência em nuvem.
\end{itemize}

%
\section{Trabalho Relacionado}
\label{sec:related}

\subsubsection{Aprendizagem por Reforço para Acompanhamento em Saúde.}
A AR tem sido aplicada a intervenções de saúde adaptativas no âmbito do paradigma
\emph{Just-In-Time Adaptive Intervention} (JITAI)~\cite{tewari2017}.
Liao et al.\ treinaram um \emph{contextual bandit} para decidir quando enviar
notificações de atividade, reduzindo o desengajamento dos
utilizadores~\cite{liao2020}. A AR profunda foi utilizada para personalizar a
prescrição de exercício em reabilitação cardíaca~\cite{escolar2022} e para
controlar protocolos de fisioterapia~\cite{yu2019}. O nosso ambiente estende estas
ideias ao modelar explicitamente a fadiga de alarmes como componente de estado de
primeira classe e ao impor restrições de segurança fortes através de penalizações
terminais.

\subsubsection{Reconhecimento de Atividade e Sensores Vestíveis.}
O Reconhecimento de Atividade Humana (HAR) a partir de unidades de medição inercial
é um problema bem estudado~\cite{banos2014}. Os smartphones modernos integram
acelerómetros, giroscópios e sensores óticos de frequência cardíaca cujos dados
podem ser fundidos para estimar a intensidade do movimento. O Dance4Life utiliza
as magnitudes brutas do acelerómetro e do giroscópio, juntamente com a frequência
cardíaca, para derivar uma pontuação de \emph{ritmo} que alimenta tanto o estado
de AR como o módulo de agrupamento de utilizadores.

\subsubsection{Sistemas Multi-Agente em Cuidados de Saúde.}
As frameworks JADE e SPADE têm sido utilizadas para construir plataformas de saúde
distribuídas, onde agentes especializados tratam da recolha de dados, raciocínio e
entrega de notificações~\cite{mea2007}. A nossa arquitetura adota o modelo de
comportamento assíncrono do SPADE, atribuindo cada preocupação lógica a um agente
dedicado.

%
\section{Formulação do Problema}
\label{sec:formulation}

Modelamos o problema de acompanhamento como um Processo de Decisão de Markov
$\mathcal{M} = (\mathcal{S}, \mathcal{A}, P, R, \gamma)$.

\subsubsection{Espaço de Estados.}
$\mathcal{S} = [0, 10]^3$ é um espaço contínuo tridimensional:
\begin{equation}
  s_t = (\text{n\'ivel\_atividade}_t,\ \text{fadiga\_f\'isica}_t,\
         \text{n\'ivel\_irrita\c{c}\~ao}_t).
\end{equation}
Todas as dimensões partilham o mesmo intervalo limitado, de modo a que a rede de
política receba entradas comensuráveis sem normalização adicional.

\subsubsection{Espaço de Ações.}
$\mathcal{A} = \{0, 1, 2, 3\}$ é um conjunto discreto de intensidades de
acompanhamento: silêncio (0), baixa (1), média (2) e alta (3).

\subsubsection{Dinâmica de Transição.}
Cada ação produz deltas determinísticos nas três dimensões do estado, acrescidos de
ruído gaussiano $\epsilon \sim \mathcal{N}(0, \sigma^2)$ para representar a
variabilidade comportamental. Ações de maior intensidade aumentam tanto o nível de
atividade como a fadiga. Um termo de acoplamento fisiológico modela o facto de que
alta atividade acelera a acumulação de fadiga. Um bónus de recuperação reduz a
fadiga quando a atividade é baixa. A Tabela~\ref{tab:dynamics} resume os deltas
nominais.

\begin{table}[h]
\caption{Deltas nominais do estado por ação (antes do ruído).}
\label{tab:dynamics}
\centering
\begin{tabular}{lccc}
\toprule
Ação & $\Delta$\,atividade & $\Delta$\,fadiga & $\Delta$\,irritação \\
\midrule
0 -- Silêncio        & $-0{,}90$ & $-0{,}80$ & $-1{,}00$ \\
1 -- Intensidade baixa  & $+0{,}40$ & $-0{,}30$ & $+0{,}20$ \\
2 -- Intensidade média  & $+1{,}20$ & $+0{,}90$ & $+0{,}70$ \\
3 -- Intensidade alta   & $+2{,}00$ & $+1{,}90$ & $+1{,}40$ \\
\bottomrule
\end{tabular}
\end{table}

\subsubsection{Função de Recompensa.}
\begin{equation}
  R(s_t, \text{terminado}) =
  \begin{cases}
    -50 & \text{se terminado (fadiga} \ge 10\ \text{ou irritação} \ge 10\text{)} \\
    +15 & \text{se}\ s_t \in \mathcal{S}_{\text{saudável}} \\
    -5  & \text{se atividade} < 2 \\
    0   & \text{caso contrário}
  \end{cases}
\end{equation}
onde $\mathcal{S}_{\text{saudável}} = \{s : \text{atividade} \in [4,7],\
\text{fadiga} < 7,\ \text{irritação} < 5\}$.

A grande penalização terminal ($-50$) encoraja a política a evitar inteiramente os
estados críticos. O bónus de zona saudável ($+15$) fornece um sinal positivo denso.
Um episódio dura no máximo 96 passos, representando um dia simulado de pontos de
decisão de acompanhamento.

%
\section{Metodologia}
\label{sec:method}

\subsection{Algoritmos}

\subsubsection{Proximal Policy Optimisation (PPO).}
O PPO~\cite{schulman2017} é um algoritmo ator-crítico \emph{on-policy} que
restringe as atualizações de política através de um objetivo substituto com
recorte (\emph{clipping}), proporcionando um treino estável. Utilizamos $n=8$
ambientes paralelos para reduzir a variância, com 1\,024 passos por \emph{rollout},
tamanho de mini-lote de 64 e 10 épocas por atualização.

\subsubsection{Deep Q-Network (DQN).}
O DQN~\cite{mnih2015} é um algoritmo baseado em valor \emph{off-policy} que aprende
a função ação-valor $Q(s,a)$ usando um \emph{replay buffer} e uma rede-alvo. O
treino utiliza um \emph{buffer} de 100\,000 transições com exploração
$\epsilon$-greedy reduzida de 1,0 para 0,05 ao longo de 25\% dos passos totais.

Ambos os algoritmos utilizam uma política de perceptrão multicamada (MLP) com duas
camadas ocultas de 64 unidades, o otimizador Adam com taxa de aprendizagem
$3 \times 10^{-4}$ (PPO) e $5 \times 10^{-4}$ (DQN), fator de desconto
$\gamma = 0{,}99$ e a mesma semente aleatória (42) para reprodutibilidade. O treino
decorre durante 400\,000 passos de ambiente em ambos os casos, utilizando a
biblioteca Stable-Baselines3~\cite{raffin2021}.

\subsection{Métricas de Avaliação}

Os modelos são avaliados a cada 10\,000 passos ao longo de 20 episódios
determinísticos. São reportadas três métricas principais:
\begin{itemize}
  \item \textbf{Recompensa média por episódio}: recompensa total não descontada,
        em média sobre os episódios de avaliação.
  \item \textbf{Duração média do episódio}: número médio de passos antes da
        terminação ou truncagem; o máximo é 96.
  \item \textbf{Taxa de sobrevivência}: fração dos episódios de avaliação que
        completaram os 96 passos sem acionar a condição de terminação crítica.
\end{itemize}

\subsection{Exportação ONNX e Implantação no Dispositivo}

Após o treino, o melhor \emph{checkpoint} PPO é exportado para o formato ONNX
usando a API \texttt{torch.onnx.export} do PyTorch. Apenas a rede de política
(ator) é exportada; a forma de entrada é $[1, 3]$, correspondendo ao vetor de
estado. O ficheiro ONNX é copiado para a diretoria \texttt{assets/models/} da
aplicação Android e carregado em tempo de execução via ONNX Runtime para Android.

%
\section{Arquitetura do Sistema}
\label{sec:arch}

O Dance4Life é composto por três camadas com acoplamento fraco: a camada de
sensores Android, o backend Python e a camada de armazenamento em nuvem (Firebase).

\subsection{Camada Android}

A aplicação Android (Kotlin, API 35) corre num smartphone, opcionalmente emparelhado
com um dispositivo Wear OS. A classe \texttt{DanceController} orquestra o
\emph{pipeline} de sensores:

\begin{enumerate}
  \item O \texttt{SensorHelper} amostra o acelerómetro, o giroscópio e o sensor
        de frequência cardíaca à taxa normal e calcula as magnitudes dos vetores.
  \item O \texttt{RitmoCalculator} mantém uma janela deslizante de 1 minuto de
        leituras de sensores e produz uma pontuação escalar de \emph{ritmo}
        representando a intensidade do movimento.
  \item O \texttt{OnnxRlCoachPolicy} constrói um tuplo
        \texttt{MovementObservation} $(a, f, i)$ a partir das métricas derivadas
        dos sensores e invoca a sessão ONNX Runtime para obter a ação de
        acompanhamento discreta.
  \item A recomendação é apresentada ao utilizador e o evento é enviado para o
        backend Python.
\end{enumerate}

O \texttt{OnnxRlCoachPolicy} mapeia a saída bruta do ONNX (\emph{logits} de ação)
para um \texttt{MovementRecommendation} contendo um título legível, duração
sugerida e uma mensagem de incentivo. O \texttt{RlMetricsManager} regista cada
evento de inferência numa fila local e envia-o para o servidor, permitindo a
melhoria futura da política de forma \emph{offline}.

\subsection{Backend Multi-Agente Python}

O backend é implementado com a \emph{framework} de agentes SPADE e uma API REST
Flask, e corre como serviço Python autónomo. A Figura~\ref{fig:arch} ilustra o
\emph{pipeline} de agentes.

\begin{figure}[h]
\centering
\begin{verbatim}
Aplicação Android
    |  POST /collect_activity
    v
[SensorAgent]
    |  SENSOR_ACTIVITY (REQUEST)
    v
[CoordinatorAgent]
    |  SENSOR_ACTIVITY (REQUEST)
    v
[HarAgent] --> [EnvironmentAgent] --> [DatabaseAgent]
               (enriquec. ambiental)    (Firebase)

[ClusteringAgent]  <-- recebe eventos de matching
    |  calcula progressão de cluster
    v
POST /set_user_match/{userId}  --> polling Android
\end{verbatim}
\caption{Fluxo de mensagens entre agentes no backend do Dance4Life.}
\label{fig:arch}
\end{figure}

\subsubsection{SensorAgent} recebe os \emph{payloads} brutos de sensores da
aplicação Android via API REST e reencaminha-os para o \texttt{CoordinatorAgent}.

\subsubsection{CoordinatorAgent} atua como hub de roteamento: despacha eventos de
sensores para o \texttt{HarAgent} para reconhecimento de atividade e para o
\texttt{ClusteringAgent} para agrupamento de utilizadores.

\subsubsection{HarAgent} reencaminha os dados de atividade para o
\texttt{EnvironmentAgent}, que enriquece o \emph{payload} com informação contextual
(temperatura meteorológica a partir de uma API meteorológica aberta) antes de o
persistir.

\subsubsection{DatabaseAgent} escreve o \emph{payload} enriquecido na Firebase
Realtime Database, fornecendo armazenamento durável e habilitando o painel web.

\subsubsection{ClusteringAgent} mantém um registo de utilizadores em memória.
Quando chega um evento de sensor, utiliza um algoritmo inspirado em k-médias para
atribuir utilizadores a um de quatro agrupamentos de atividade---\emph{Iniciante},
\emph{Moderado}, \emph{Avançado}, \emph{Expert}---com base na sua pontuação de
\emph{ritmo} acumulada. O agente gera então um convite de progressão de
agrupamento, orientando o utilizador para o nível seguinte, enviado de volta ao
cliente Android via \texttt{/set\_user\_match}.

\subsection{Painel Web}

Uma aplicação web estática separada (HTML/CSS/JavaScript) liga-se diretamente ao
Firebase e apresenta visualizações em tempo real dos dados de atividade, resultados
de agrupamento de utilizadores e recomendações de movimento, destinada a
cuidadores e coordenadores de programas.

%
\section{Resultados}
\label{sec:results}

\subsection{Convergência do Treino}

A Tabela~\ref{tab:results} resume os resultados de avaliação recolhidos a cada
10\,000 passos. Ambos os algoritmos convergem rapidamente. O DQN atinge uma
recompensa média acima de 1\,430 logo aos 10\,000 passos, devido à sua natureza
\emph{off-policy} e ao espaço de estados relativamente simples. O PPO demora
ligeiramente mais a estabilizar---atingindo a mesma banda de recompensa
aproximadamente aos 50\,000 passos---mas apresenta menor variância no regime
convergido.

\begin{table}[h]
\caption{Métricas de avaliação finais (média sobre os últimos 10
\emph{checkpoints} de avaliação, após convergência).}
\label{tab:results}
\centering
\begin{tabular}{lcccc}
\toprule
Algoritmo & Recomp.\ Média & Desvio Padrão & Dur.\ Ep. & Taxa Sobrev. \\
\midrule
PPO & $1\,434{,}5$ & $\pm 2{,}1$ & 96 & 100\% \\
DQN & $1\,432{,}8$ & $\pm 4{,}3$ & 96 & 100\% \\
\bottomrule
\end{tabular}
\end{table}

Ambas as políticas atingem uma taxa de sobrevivência de 100\%: nenhum episódio
termina antes do limite de 96 passos, o que significa que nem a fadiga crítica
nem a irritação crítica são alguma vez acionadas. A duração média dos episódios é
identicamente 96 para todos os \emph{checkpoints} avaliados após a convergência,
confirmando o domínio completo da restrição de segurança.

A recompensa máxima teórica para um episódio em que o agente permanece na zona
saudável durante todos os 96 passos é $96 \times 15 = 1\,440$. Ambos os algoritmos
atingem aproximadamente 99,7\% deste limite, indicando que as políticas aprendidas
mantêm o utilizador na zona saudável de forma quase contínua.

\subsection{Distribuição de Ações}

Ambas as políticas selecionam predominantemente o acompanhamento de
\emph{intensidade média} (ação 2), usando ocasionalmente o \emph{silêncio}
(ação 0) para prevenir a acumulação de irritação e a \emph{intensidade baixa}
(ação 1) durante fases de recuperação. A \emph{intensidade alta} (ação 3) é
utilizada de forma parcimoniosa, pois os seus grandes deltas de fadiga e irritação
arriscam a penalização terminal.

\subsection{Validação da Implantação}

O modelo PPO ONNX exportado (\texttt{dance4life\_coach\_v2\_ppo.onnx}) foi validado
comparando o argmax dos \emph{logits} ONNX com a saída do método \texttt{predict}
do Stable-Baselines3 para 100 observações aleatórias. Não foi encontrada nenhuma
discrepância, confirmando uma exportação sem perdas. O tamanho do modelo é de
aproximadamente 18\,KB, adequado para inferência no dispositivo com latência
negligenciável em hardware Android de gama média.

\subsection{Discussão}

Os resultados demonstram que a formulação de AR é bem adequada a este problema de
acompanhamento. A recompensa de zona saudável fornece um sinal de aprendizagem
denso que orienta rapidamente a política para a região de operação desejada. A
grande penalização terminal ensina eficazmente a política a tratar a fadiga de
alarmes e a sobrecarga física como restrições rígidas, em vez de compromissos
suaves.

O desempenho quase idêntico de PPO e DQN é explicado pelo espaço de estados
pequeno e contínuo e pela ausência de observabilidade parcial: o estado determina
completamente a ação ótima. Em implantações mais realistas, onde a observação é
ruidosa ou parcialmente observada, a exploração mais forte do PPO através da
regularização por entropia pode conferir uma vantagem.

Uma limitação do ambiente atual é o facto de o modelo de transição de estados ser
construído manualmente. Calibrar as magnitudes de ruído e os coeficientes de
acoplamento com dados fisiológicos reais de utilizadores aumentaria a validade
ecológica. Além disso, a observação fornecida à política Android é derivada de
sinais de sensores através de mapeamentos heurísticos; trabalho futuro poderia
aprender este mapeamento de ponta a ponta a partir de dados fisiológicos rotulados.

%
\section{Conclusões}
\label{sec:conclusions}

Apresentámos o Dance4Life, uma plataforma de ponta a ponta que aplica a
Aprendizagem por Reforço à personalização do acompanhamento de atividade física
para o envelhecimento ativo através da dança. Um ambiente Gymnasium personalizado
codifica o compromisso a três vias entre promoção da atividade, segurança física
e envolvimento do utilizador. Tanto o PPO como o DQN convergem para um desempenho
próximo do ótimo, atingindo uma taxa de sobrevivência de 100\% e aproximadamente
99,7\% da recompensa máxima alcançável. A política PPO aprendida é exportada para
ONNX e implantada no Android via ONNX Runtime, permitindo acompanhamento em tempo
real no dispositivo a partir de dados de sensores vestíveis. Um backend multi-agente
SPADE trata do restante pipeline: roteamento baseado em agentes, enriquecimento de
contexto ambiental, agrupamento de utilizadores por k-médias com acompanhamento de
nível progressivo e persistência em nuvem Firebase.

O trabalho futuro incidirá sobre: (i)~recolha de dados reais de utilizadores para
calibrar a dinâmica do ambiente; (ii)~substituição do mapeamento heurístico de
estado por um codificador sensor-estado aprendido; (iii)~avaliação do sistema
completo com uma coorte de voluntários adultos mais velhos; e (iv)~extensão do
módulo Wear OS para fornecer \emph{feedback} de frequência cardíaca mais granular
ao estado de AR.

%
\begin{credits}
\subsubsection{\discintname}
Os autores não têm interesses concorrentes a declarar.
\end{credits}

%
% ---- Bibliography ----
%
\begin{thebibliography}{10}

\bibitem{who2022}
World Health Organization: Ageing and Health. WHO Fact Sheet (2022).
\url{https://www.who.int/news-room/fact-sheets/detail/ageing-and-health}

\bibitem{hwang2016}
Hwang, P.W.-N., Braun, K.L.: The effectiveness of dance interventions to
improve older adults' health: a systematic literature review. Altern.\ Ther.\
Health Med.\ \textbf{21}(5), 64--70 (2015)

\bibitem{tewari2017}
Tewari, A., Murphy, S.A.: From user engagement to health behaviour: just-in-time
adaptive interventions. In: mHealth, pp. 167--196. Springer (2017)

\bibitem{liao2020}
Liao, P., et al.: Personalized HeartSteps: a reinforcement learning algorithm
for optimizing physical activity. Proc.\ ACM Interact.\ Mob.\ Wearable Ubiquitous
Technol.\ \textbf{4}(1), 18:1--18:22 (2020)

\bibitem{escolar2022}
Escolar-Reina, P., et al.: Artificial intelligence for exercise prescription in
rehabilitation: challenges and perspectives. Front.\ Physiol.\ \textbf{13},
900697 (2022)

\bibitem{yu2019}
Yu, C., et al.: Reinforcement learning in healthcare: a survey. ACM Comput.\ Surv.\
\textbf{55}(1), 1--36 (2023)

\bibitem{banos2014}
Banos, O., et al.: Window size impact in human activity recognition. Sensors
\textbf{14}(4), 6474--6499 (2014)

\bibitem{mea2007}
Mea, V.D.: Agents in medicine. Stud.\ Health Technol.\ Inform.\ \textbf{124},
791--796 (2006)

\bibitem{schulman2017}
Schulman, J., et al.: Proximal policy optimization algorithms.
arXiv:1707.06347 (2017)

\bibitem{mnih2015}
Mnih, V., et al.: Human-level control through deep reinforcement learning.
Nature \textbf{518}(7540), 529--533 (2015)

\bibitem{raffin2021}
Raffin, A., et al.: Stable-Baselines3: reliable reinforcement learning
implementations. J.\ Mach.\ Learn.\ Res.\ \textbf{22}(268), 1--8 (2021)

\end{thebibliography}
\end{document}
